% paper79-prompt-to-proof.tex — From Prompt to Proof: Compiling Natural
%   Language Intent into Sheaf Certificates
%   Paper 79 of the Judgment Geometry series.
% Compile: pdflatex paper79-prompt-to-proof && pdflatex paper79-prompt-to-proof
\documentclass[11pt]{article}
\usepackage{lmodern}
% jugeo-common.tex — shared preamble for all Judgment Geometry papers
% Include via: % jugeo-common.tex — shared preamble for all Judgment Geometry papers
% Include via: % jugeo-common.tex — shared preamble for all Judgment Geometry papers
% Include via: \input{jugeo-common}

% ── Packages ────────────────────────────────────────────────────────────
\usepackage{amsmath,amssymb,amsthm,mathtools}
\usepackage{tikz-cd}
\usepackage{listings}
\usepackage{xcolor}
\usepackage{xspace}
\usepackage{booktabs}
\usepackage{multirow}
\usepackage{enumitem}
\usepackage[expansion=false]{microtype}
\usepackage{hyperref}
\usepackage{cleveref}
% \usepackage{thmtools}  % disabled: not available in this TeX installation
\usepackage{float}
\usepackage{caption}
\usepackage{subcaption}
\usepackage{tabularx}
\usepackage{stmaryrd}       % \llbracket, \rrbracket
\usepackage{bbm}            % \mathbbm{1}

% ── Page geometry ───────────────────────────────────────────────────────
\usepackage[margin=1in]{geometry}

% ── Theorem environments ────────────────────────────────────────────────
\theoremstyle{plain}
\newtheorem{theorem}{Theorem}[section]
\newtheorem{lemma}[theorem]{Lemma}
\newtheorem{proposition}[theorem]{Proposition}
\newtheorem{corollary}[theorem]{Corollary}

\theoremstyle{definition}
\newtheorem{definition}[theorem]{Definition}
\newtheorem{example}[theorem]{Example}
\newtheorem{construction}[theorem]{Construction}

\theoremstyle{remark}
\newtheorem{remark}[theorem]{Remark}
\newtheorem{notation}[theorem]{Notation}

% ── Python listings style ──────────────────────────────────────────────
\definecolor{jugeoBlue}{HTML}{2563EB}
\definecolor{jugeoGreen}{HTML}{059669}
\definecolor{jugeoRed}{HTML}{DC2626}
\definecolor{jugeoPurple}{HTML}{7C3AED}
\definecolor{jugeoGray}{HTML}{6B7280}
\definecolor{jugeoBg}{HTML}{F8FAFC}

\lstdefinestyle{jugeo-python}{
  language=Python,
  basicstyle=\ttfamily\small,
  keywordstyle=\color{jugeoBlue}\bfseries,
  stringstyle=\color{jugeoGreen},
  commentstyle=\color{jugeoGray}\itshape,
  numberstyle=\tiny\color{jugeoGray},
  backgroundcolor=\color{jugeoBg},
  numbers=left,
  numbersep=6pt,
  frame=single,
  framerule=0.4pt,
  rulecolor=\color{jugeoGray!40},
  breaklines=true,
  breakatwhitespace=true,
  showstringspaces=false,
  tabsize=4,
  xleftmargin=1.5em,
  framexleftmargin=1.5em,
  morekeywords={dataclass,frozen,field,Enum,IntEnum,auto,Any,
                Mapping,Sequence,Optional,match,case},
  literate={→}{{$\to$}}1 {←}{{$\leftarrow$}}1
           {≤}{{$\leq$}}1 {≥}{{$\geq$}}1 {≠}{{$\neq$}}1
           {∈}{{$\in$}}1 {∀}{{$\forall$}}1 {∃}{{$\exists$}}1
           {⊕}{{$\oplus$}}1 {⊖}{{$\ominus$}}1,
  escapeinside={(*@}{@*)},
}
\lstset{style=jugeo-python}

\lstdefinestyle{jugeo-output}{
  basicstyle=\ttfamily\small,
  backgroundcolor=\color{jugeoBg},
  frame=single,
  framerule=0.4pt,
  rulecolor=\color{jugeoGray!40},
  breaklines=true,
  showstringspaces=false,
  xleftmargin=1.5em,
  framexleftmargin=0.5em,
  numbers=none,
}

% ── JuGeo-specific macros ──────────────────────────────────────────────

% Judgment 8-tuple
\newcommand{\J}{\mathcal{J}}
\newcommand{\Jud}[8]{(#1,\,#2,\,#3,\,#4,\,#5,\,#6,\,#7,\,#8)}
\newcommand{\JudShort}[1]{\J_{#1}}

% Components
\newcommand{\coord}{\mathsf{c}}            % coordinate
\newcommand{\prop}{\varphi}                 % proposition
\newcommand{\carrier}{\mathcal{A}}          % carrier type
\newcommand{\evid}{\mathcal{E}}             % evidence bundle
\newcommand{\oblig}{\mathcal{O}}            % obligations
\newcommand{\obstr}{\mathcal{B}}            % obstructions
\newcommand{\trust}{\mathcal{T}}            % trust annotation
\newcommand{\prov}{\Pi}                     % provenance

% Trust algebra
\newcommand{\Talg}{\mathfrak{T}}            % trust ordered algebra
\newcommand{\Eadm}{\mathcal{E}_{\mathrm{adm}}}
\newcommand{\tleq}{\preceq}                 % trust ordering
\newcommand{\tjoin}{\oplus}                 % conservative join
\newcommand{\tatten}{\ominus}               % attenuation
\newcommand{\tpromote}{\uparrow_\pi}        % promotion
\newcommand{\tdemote}{\downarrow_\chi}       % demotion

% Trust tiers
\newcommand{\Tcontra}{\ensuremath{\mathsf{CONTRADICTED}}}
\newcommand{\Tunver}{\ensuremath{\mathsf{UNVERIFIED}}}
\newcommand{\Tcopilot}{\ensuremath{\mathsf{COPILOT}}}
\newcommand{\Toracle}{\ensuremath{\mathsf{ORACLE}}}
\newcommand{\Truntime}{\ensuremath{\mathsf{RUNTIME}}}
\newcommand{\Tsolver}{\ensuremath{\mathsf{SOLVER}}}
\newcommand{\Tproof}{\ensuremath{\mathsf{PROOF}}}

% Site and geometry
\newcommand{\Site}{\mathcal{S}}             % semantic site
\newcommand{\Cov}{\mathrm{Cov}}            % covering family
\newcommand{\Ob}{\mathrm{Ob}}              % objects
\newcommand{\Mor}{\mathrm{Mor}}            % morphisms
\newcommand{\res}{\mathrm{res}}            % restriction
\newcommand{\Cech}{\check{C}}              % Čech complex
\newcommand{\Hcoh}[1]{H^{#1}}             % cohomology group

% Descent
\newcommand{\descent}{\mathrm{desc}}
\newcommand{\glue}{\mathrm{glue}}
\newcommand{\overlap}[2]{#1 \cap #2}

% Semantic moves
\newcommand{\VERIFY}{\textsc{Verify}}
\newcommand{\CONSTRUCT}{\textsc{Construct}}
\newcommand{\NORMALIZE}{\textsc{Normalize}}
\newcommand{\GLUE}{\textsc{Glue}}
\newcommand{\RESTRICT}{\textsc{Restrict}}
\newcommand{\TRANSPORT}{\textsc{Transport}}
\newcommand{\REPAIR}{\textsc{Repair}}
\newcommand{\PROMOTE}{\textsc{Promote}}
\newcommand{\CHALLENGE}{\textsc{Challenge}}
\newcommand{\ESCALATE}{\textsc{Escalate}}

% Fragments
\newcommand{\QFLIA}{\texttt{QF\_LIA}}
\newcommand{\QFLRA}{\texttt{QF\_LRA}}
\newcommand{\QFBV}{\texttt{QF\_BV}}
\newcommand{\QFUF}{\texttt{QF\_UF}}

% Misc
\newcommand{\Python}{\textsc{Python}}
\newcommand{\JuGeo}{\textsc{JuGeo}}
\newcommand{\LEAN}{\textsc{Lean}\,4}
\newcommand{\Fstar}{F$^{\star}$}
\newcommand{\Coq}{\textsc{Coq}}
\newcommand{\Dafny}{\textsc{Dafny}}

% Common shortcuts
\newcommand{\ie}{i.e.\@\xspace}
\newcommand{\eg}{e.g.\@\xspace}
\newcommand{\cf}{cf.\@\xspace}
\newcommand{\wrt}{w.r.t.\@\xspace}

% Evaluation shortcuts
\newcommand{\numcases}{300}
\newcommand{\numspec}{100}
\newcommand{\numeq}{100}
\newcommand{\numbug}{100}
\newcommand{\medtime}{6.5\,ms}
\newcommand{\totaltime}{1.949\,s}


% ── Packages ────────────────────────────────────────────────────────────
\usepackage{amsmath,amssymb,amsthm,mathtools}
\usepackage{tikz-cd}
\usepackage{listings}
\usepackage{xcolor}
\usepackage{xspace}
\usepackage{booktabs}
\usepackage{multirow}
\usepackage{enumitem}
\usepackage[expansion=false]{microtype}
\usepackage{hyperref}
\usepackage{cleveref}
% \usepackage{thmtools}  % disabled: not available in this TeX installation
\usepackage{float}
\usepackage{caption}
\usepackage{subcaption}
\usepackage{tabularx}
\usepackage{stmaryrd}       % \llbracket, \rrbracket
\usepackage{bbm}            % \mathbbm{1}

% ── Page geometry ───────────────────────────────────────────────────────
\usepackage[margin=1in]{geometry}

% ── Theorem environments ────────────────────────────────────────────────
\theoremstyle{plain}
\newtheorem{theorem}{Theorem}[section]
\newtheorem{lemma}[theorem]{Lemma}
\newtheorem{proposition}[theorem]{Proposition}
\newtheorem{corollary}[theorem]{Corollary}

\theoremstyle{definition}
\newtheorem{definition}[theorem]{Definition}
\newtheorem{example}[theorem]{Example}
\newtheorem{construction}[theorem]{Construction}

\theoremstyle{remark}
\newtheorem{remark}[theorem]{Remark}
\newtheorem{notation}[theorem]{Notation}

% ── Python listings style ──────────────────────────────────────────────
\definecolor{jugeoBlue}{HTML}{2563EB}
\definecolor{jugeoGreen}{HTML}{059669}
\definecolor{jugeoRed}{HTML}{DC2626}
\definecolor{jugeoPurple}{HTML}{7C3AED}
\definecolor{jugeoGray}{HTML}{6B7280}
\definecolor{jugeoBg}{HTML}{F8FAFC}

\lstdefinestyle{jugeo-python}{
  language=Python,
  basicstyle=\ttfamily\small,
  keywordstyle=\color{jugeoBlue}\bfseries,
  stringstyle=\color{jugeoGreen},
  commentstyle=\color{jugeoGray}\itshape,
  numberstyle=\tiny\color{jugeoGray},
  backgroundcolor=\color{jugeoBg},
  numbers=left,
  numbersep=6pt,
  frame=single,
  framerule=0.4pt,
  rulecolor=\color{jugeoGray!40},
  breaklines=true,
  breakatwhitespace=true,
  showstringspaces=false,
  tabsize=4,
  xleftmargin=1.5em,
  framexleftmargin=1.5em,
  morekeywords={dataclass,frozen,field,Enum,IntEnum,auto,Any,
                Mapping,Sequence,Optional,match,case},
  literate={→}{{$\to$}}1 {←}{{$\leftarrow$}}1
           {≤}{{$\leq$}}1 {≥}{{$\geq$}}1 {≠}{{$\neq$}}1
           {∈}{{$\in$}}1 {∀}{{$\forall$}}1 {∃}{{$\exists$}}1
           {⊕}{{$\oplus$}}1 {⊖}{{$\ominus$}}1,
  escapeinside={(*@}{@*)},
}
\lstset{style=jugeo-python}

\lstdefinestyle{jugeo-output}{
  basicstyle=\ttfamily\small,
  backgroundcolor=\color{jugeoBg},
  frame=single,
  framerule=0.4pt,
  rulecolor=\color{jugeoGray!40},
  breaklines=true,
  showstringspaces=false,
  xleftmargin=1.5em,
  framexleftmargin=0.5em,
  numbers=none,
}

% ── JuGeo-specific macros ──────────────────────────────────────────────

% Judgment 8-tuple
\newcommand{\J}{\mathcal{J}}
\newcommand{\Jud}[8]{(#1,\,#2,\,#3,\,#4,\,#5,\,#6,\,#7,\,#8)}
\newcommand{\JudShort}[1]{\J_{#1}}

% Components
\newcommand{\coord}{\mathsf{c}}            % coordinate
\newcommand{\prop}{\varphi}                 % proposition
\newcommand{\carrier}{\mathcal{A}}          % carrier type
\newcommand{\evid}{\mathcal{E}}             % evidence bundle
\newcommand{\oblig}{\mathcal{O}}            % obligations
\newcommand{\obstr}{\mathcal{B}}            % obstructions
\newcommand{\trust}{\mathcal{T}}            % trust annotation
\newcommand{\prov}{\Pi}                     % provenance

% Trust algebra
\newcommand{\Talg}{\mathfrak{T}}            % trust ordered algebra
\newcommand{\Eadm}{\mathcal{E}_{\mathrm{adm}}}
\newcommand{\tleq}{\preceq}                 % trust ordering
\newcommand{\tjoin}{\oplus}                 % conservative join
\newcommand{\tatten}{\ominus}               % attenuation
\newcommand{\tpromote}{\uparrow_\pi}        % promotion
\newcommand{\tdemote}{\downarrow_\chi}       % demotion

% Trust tiers
\newcommand{\Tcontra}{\ensuremath{\mathsf{CONTRADICTED}}}
\newcommand{\Tunver}{\ensuremath{\mathsf{UNVERIFIED}}}
\newcommand{\Tcopilot}{\ensuremath{\mathsf{COPILOT}}}
\newcommand{\Toracle}{\ensuremath{\mathsf{ORACLE}}}
\newcommand{\Truntime}{\ensuremath{\mathsf{RUNTIME}}}
\newcommand{\Tsolver}{\ensuremath{\mathsf{SOLVER}}}
\newcommand{\Tproof}{\ensuremath{\mathsf{PROOF}}}

% Site and geometry
\newcommand{\Site}{\mathcal{S}}             % semantic site
\newcommand{\Cov}{\mathrm{Cov}}            % covering family
\newcommand{\Ob}{\mathrm{Ob}}              % objects
\newcommand{\Mor}{\mathrm{Mor}}            % morphisms
\newcommand{\res}{\mathrm{res}}            % restriction
\newcommand{\Cech}{\check{C}}              % Čech complex
\newcommand{\Hcoh}[1]{H^{#1}}             % cohomology group

% Descent
\newcommand{\descent}{\mathrm{desc}}
\newcommand{\glue}{\mathrm{glue}}
\newcommand{\overlap}[2]{#1 \cap #2}

% Semantic moves
\newcommand{\VERIFY}{\textsc{Verify}}
\newcommand{\CONSTRUCT}{\textsc{Construct}}
\newcommand{\NORMALIZE}{\textsc{Normalize}}
\newcommand{\GLUE}{\textsc{Glue}}
\newcommand{\RESTRICT}{\textsc{Restrict}}
\newcommand{\TRANSPORT}{\textsc{Transport}}
\newcommand{\REPAIR}{\textsc{Repair}}
\newcommand{\PROMOTE}{\textsc{Promote}}
\newcommand{\CHALLENGE}{\textsc{Challenge}}
\newcommand{\ESCALATE}{\textsc{Escalate}}

% Fragments
\newcommand{\QFLIA}{\texttt{QF\_LIA}}
\newcommand{\QFLRA}{\texttt{QF\_LRA}}
\newcommand{\QFBV}{\texttt{QF\_BV}}
\newcommand{\QFUF}{\texttt{QF\_UF}}

% Misc
\newcommand{\Python}{\textsc{Python}}
\newcommand{\JuGeo}{\textsc{JuGeo}}
\newcommand{\LEAN}{\textsc{Lean}\,4}
\newcommand{\Fstar}{F$^{\star}$}
\newcommand{\Coq}{\textsc{Coq}}
\newcommand{\Dafny}{\textsc{Dafny}}

% Common shortcuts
\newcommand{\ie}{i.e.\@\xspace}
\newcommand{\eg}{e.g.\@\xspace}
\newcommand{\cf}{cf.\@\xspace}
\newcommand{\wrt}{w.r.t.\@\xspace}

% Evaluation shortcuts
\newcommand{\numcases}{300}
\newcommand{\numspec}{100}
\newcommand{\numeq}{100}
\newcommand{\numbug}{100}
\newcommand{\medtime}{6.5\,ms}
\newcommand{\totaltime}{1.949\,s}


% ── Packages ────────────────────────────────────────────────────────────
\usepackage{amsmath,amssymb,amsthm,mathtools}
\usepackage{tikz-cd}
\usepackage{listings}
\usepackage{xcolor}
\usepackage{xspace}
\usepackage{booktabs}
\usepackage{multirow}
\usepackage{enumitem}
\usepackage[expansion=false]{microtype}
\usepackage{hyperref}
\usepackage{cleveref}
% \usepackage{thmtools}  % disabled: not available in this TeX installation
\usepackage{float}
\usepackage{caption}
\usepackage{subcaption}
\usepackage{tabularx}
\usepackage{stmaryrd}       % \llbracket, \rrbracket
\usepackage{bbm}            % \mathbbm{1}

% ── Page geometry ───────────────────────────────────────────────────────
\usepackage[margin=1in]{geometry}

% ── Theorem environments ────────────────────────────────────────────────
\theoremstyle{plain}
\newtheorem{theorem}{Theorem}[section]
\newtheorem{lemma}[theorem]{Lemma}
\newtheorem{proposition}[theorem]{Proposition}
\newtheorem{corollary}[theorem]{Corollary}

\theoremstyle{definition}
\newtheorem{definition}[theorem]{Definition}
\newtheorem{example}[theorem]{Example}
\newtheorem{construction}[theorem]{Construction}

\theoremstyle{remark}
\newtheorem{remark}[theorem]{Remark}
\newtheorem{notation}[theorem]{Notation}

% ── Python listings style ──────────────────────────────────────────────
\definecolor{jugeoBlue}{HTML}{2563EB}
\definecolor{jugeoGreen}{HTML}{059669}
\definecolor{jugeoRed}{HTML}{DC2626}
\definecolor{jugeoPurple}{HTML}{7C3AED}
\definecolor{jugeoGray}{HTML}{6B7280}
\definecolor{jugeoBg}{HTML}{F8FAFC}

\lstdefinestyle{jugeo-python}{
  language=Python,
  basicstyle=\ttfamily\small,
  keywordstyle=\color{jugeoBlue}\bfseries,
  stringstyle=\color{jugeoGreen},
  commentstyle=\color{jugeoGray}\itshape,
  numberstyle=\tiny\color{jugeoGray},
  backgroundcolor=\color{jugeoBg},
  numbers=left,
  numbersep=6pt,
  frame=single,
  framerule=0.4pt,
  rulecolor=\color{jugeoGray!40},
  breaklines=true,
  breakatwhitespace=true,
  showstringspaces=false,
  tabsize=4,
  xleftmargin=1.5em,
  framexleftmargin=1.5em,
  morekeywords={dataclass,frozen,field,Enum,IntEnum,auto,Any,
                Mapping,Sequence,Optional,match,case},
  literate={→}{{$\to$}}1 {←}{{$\leftarrow$}}1
           {≤}{{$\leq$}}1 {≥}{{$\geq$}}1 {≠}{{$\neq$}}1
           {∈}{{$\in$}}1 {∀}{{$\forall$}}1 {∃}{{$\exists$}}1
           {⊕}{{$\oplus$}}1 {⊖}{{$\ominus$}}1,
  escapeinside={(*@}{@*)},
}
\lstset{style=jugeo-python}

\lstdefinestyle{jugeo-output}{
  basicstyle=\ttfamily\small,
  backgroundcolor=\color{jugeoBg},
  frame=single,
  framerule=0.4pt,
  rulecolor=\color{jugeoGray!40},
  breaklines=true,
  showstringspaces=false,
  xleftmargin=1.5em,
  framexleftmargin=0.5em,
  numbers=none,
}

% ── JuGeo-specific macros ──────────────────────────────────────────────

% Judgment 8-tuple
\newcommand{\J}{\mathcal{J}}
\newcommand{\Jud}[8]{(#1,\,#2,\,#3,\,#4,\,#5,\,#6,\,#7,\,#8)}
\newcommand{\JudShort}[1]{\J_{#1}}

% Components
\newcommand{\coord}{\mathsf{c}}            % coordinate
\newcommand{\prop}{\varphi}                 % proposition
\newcommand{\carrier}{\mathcal{A}}          % carrier type
\newcommand{\evid}{\mathcal{E}}             % evidence bundle
\newcommand{\oblig}{\mathcal{O}}            % obligations
\newcommand{\obstr}{\mathcal{B}}            % obstructions
\newcommand{\trust}{\mathcal{T}}            % trust annotation
\newcommand{\prov}{\Pi}                     % provenance

% Trust algebra
\newcommand{\Talg}{\mathfrak{T}}            % trust ordered algebra
\newcommand{\Eadm}{\mathcal{E}_{\mathrm{adm}}}
\newcommand{\tleq}{\preceq}                 % trust ordering
\newcommand{\tjoin}{\oplus}                 % conservative join
\newcommand{\tatten}{\ominus}               % attenuation
\newcommand{\tpromote}{\uparrow_\pi}        % promotion
\newcommand{\tdemote}{\downarrow_\chi}       % demotion

% Trust tiers
\newcommand{\Tcontra}{\ensuremath{\mathsf{CONTRADICTED}}}
\newcommand{\Tunver}{\ensuremath{\mathsf{UNVERIFIED}}}
\newcommand{\Tcopilot}{\ensuremath{\mathsf{COPILOT}}}
\newcommand{\Toracle}{\ensuremath{\mathsf{ORACLE}}}
\newcommand{\Truntime}{\ensuremath{\mathsf{RUNTIME}}}
\newcommand{\Tsolver}{\ensuremath{\mathsf{SOLVER}}}
\newcommand{\Tproof}{\ensuremath{\mathsf{PROOF}}}

% Site and geometry
\newcommand{\Site}{\mathcal{S}}             % semantic site
\newcommand{\Cov}{\mathrm{Cov}}            % covering family
\newcommand{\Ob}{\mathrm{Ob}}              % objects
\newcommand{\Mor}{\mathrm{Mor}}            % morphisms
\newcommand{\res}{\mathrm{res}}            % restriction
\newcommand{\Cech}{\check{C}}              % Čech complex
\newcommand{\Hcoh}[1]{H^{#1}}             % cohomology group

% Descent
\newcommand{\descent}{\mathrm{desc}}
\newcommand{\glue}{\mathrm{glue}}
\newcommand{\overlap}[2]{#1 \cap #2}

% Semantic moves
\newcommand{\VERIFY}{\textsc{Verify}}
\newcommand{\CONSTRUCT}{\textsc{Construct}}
\newcommand{\NORMALIZE}{\textsc{Normalize}}
\newcommand{\GLUE}{\textsc{Glue}}
\newcommand{\RESTRICT}{\textsc{Restrict}}
\newcommand{\TRANSPORT}{\textsc{Transport}}
\newcommand{\REPAIR}{\textsc{Repair}}
\newcommand{\PROMOTE}{\textsc{Promote}}
\newcommand{\CHALLENGE}{\textsc{Challenge}}
\newcommand{\ESCALATE}{\textsc{Escalate}}

% Fragments
\newcommand{\QFLIA}{\texttt{QF\_LIA}}
\newcommand{\QFLRA}{\texttt{QF\_LRA}}
\newcommand{\QFBV}{\texttt{QF\_BV}}
\newcommand{\QFUF}{\texttt{QF\_UF}}

% Misc
\newcommand{\Python}{\textsc{Python}}
\newcommand{\JuGeo}{\textsc{JuGeo}}
\newcommand{\LEAN}{\textsc{Lean}\,4}
\newcommand{\Fstar}{F$^{\star}$}
\newcommand{\Coq}{\textsc{Coq}}
\newcommand{\Dafny}{\textsc{Dafny}}

% Common shortcuts
\newcommand{\ie}{i.e.\@\xspace}
\newcommand{\eg}{e.g.\@\xspace}
\newcommand{\cf}{cf.\@\xspace}
\newcommand{\wrt}{w.r.t.\@\xspace}

% Evaluation shortcuts
\newcommand{\numcases}{300}
\newcommand{\numspec}{100}
\newcommand{\numeq}{100}
\newcommand{\numbug}{100}
\newcommand{\medtime}{6.5\,ms}
\newcommand{\totaltime}{1.949\,s}

% data-paper79.tex — Experiment data for Paper 79
% Prompt to Proof: Compiling Natural Language Intent into Sheaf Certificates

% ── Overall pipeline statistics ────────────────────────────────────────
\newcommand{\ptpTotalPrompts}{400}
\newcommand{\ptpSuccessCert}{347}
\newcommand{\ptpSuccessRate}{86.8\%}
\newcommand{\ptpPartialCert}{31}
\newcommand{\ptpPartialRate}{7.8\%}
\newcommand{\ptpFailedCert}{22}
\newcommand{\ptpFailedRate}{5.5\%}

% ── Per-benchmark breakdown ────────────────────────────────────────────
\newcommand{\ptpHumanEvalTotal}{164}
\newcommand{\ptpHumanEvalCert}{148}
\newcommand{\ptpHumanEvalRate}{90.2\%}
\newcommand{\ptpHumanEvalPartial}{10}
\newcommand{\ptpHumanEvalFailed}{6}

\newcommand{\ptpMBPPTotal}{128}
\newcommand{\ptpMBPPCert}{108}
\newcommand{\ptpMBPPRate}{84.4\%}
\newcommand{\ptpMBPPPartial}{12}
\newcommand{\ptpMBPPFailed}{8}

\newcommand{\ptpCustomTotal}{108}
\newcommand{\ptpCustomCert}{91}
\newcommand{\ptpCustomRate}{84.3\%}
\newcommand{\ptpCustomPartial}{9}
\newcommand{\ptpCustomFailed}{8}

% ── Stage 1: Prompt parsing ───────────────────────────────────────────
\newcommand{\ptpParseTimeMedian}{42\,ms}
\newcommand{\ptpParseTimeMean}{58\,ms}
\newcommand{\ptpParseTimeMax}{310\,ms}
\newcommand{\ptpParseIntentsMean}{3.7}
\newcommand{\ptpParseIntentsMedian}{3}
\newcommand{\ptpParseIntentsMax}{12}
\newcommand{\ptpParseFragmentsMean}{5.1}

% ── Stage 2: LLM code generation ──────────────────────────────────────
\newcommand{\ptpGenTimeMedian}{1.8\,s}
\newcommand{\ptpGenTimeMean}{2.3\,s}
\newcommand{\ptpGenTimeMax}{8.4\,s}
\newcommand{\ptpGenPassRate}{94.5\%}
\newcommand{\ptpGenLOCMean}{28}
\newcommand{\ptpGenLOCMedian}{22}

% ── Stage 3: Site construction ────────────────────────────────────────
\newcommand{\ptpSiteTimeMedian}{320\,ms}
\newcommand{\ptpSiteTimeMean}{410\,ms}
\newcommand{\ptpSiteTimeMax}{2.1\,s}
\newcommand{\ptpSiteCoordsMean}{5.8}
\newcommand{\ptpSiteCoordsMedian}{5}
\newcommand{\ptpSiteCoordsMax}{18}
\newcommand{\ptpSitePropsMean}{11.3}
\newcommand{\ptpSitePropsMedian}{10}
\newcommand{\ptpSitePropsMax}{34}

% ── Stage 4: Descent verification ─────────────────────────────────────
\newcommand{\ptpDescentTimeMedian}{1.4\,s}
\newcommand{\ptpDescentTimeMean}{2.1\,s}
\newcommand{\ptpDescentTimeMax}{12.6\,s}
\newcommand{\ptpDescentCocycleRate}{91.2\%}
\newcommand{\ptpDescentGlueRate}{89.5\%}
\newcommand{\ptpDescentObstructions}{47}
\newcommand{\ptpDescentHOneZero}{353}

% ── Stage 5: Certificate emission ─────────────────────────────────────
\newcommand{\ptpCertTimeMedian}{85\,ms}
\newcommand{\ptpCertTimeMean}{120\,ms}
\newcommand{\ptpCertTimeMax}{480\,ms}
\newcommand{\ptpCertSizeMedianKB}{2.4\,KB}
\newcommand{\ptpCertSizeMeanKB}{3.8\,KB}
\newcommand{\ptpCertSizeMaxKB}{18.2\,KB}
\newcommand{\ptpCertOverhead}{1.6$\times$}

% ── End-to-end pipeline ──────────────────────────────────────────────
\newcommand{\ptpE2ETimeMedian}{4.2\,s}
\newcommand{\ptpE2ETimeMean}{5.6\,s}
\newcommand{\ptpE2ETimeMax}{21.8\,s}
\newcommand{\ptpE2ETimeTotal}{37.3\,min}

% ── Verification quality ─────────────────────────────────────────────
\newcommand{\ptpFalsePositive}{0.6\%}
\newcommand{\ptpFalseNegative}{3.2\%}
\newcommand{\ptpReverifMedian}{12\,\textmu{}s}
\newcommand{\ptpReverifMean}{18\,\textmu{}s}
\newcommand{\ptpReverifMax}{95\,\textmu{}s}

% ── Certificate composition ──────────────────────────────────────────
\newcommand{\ptpComposeTestTotal}{80}
\newcommand{\ptpComposeSuccRate}{92.5\%}
\newcommand{\ptpComposeTimeMean}{180\,ms}

% ── Trust tier distribution ──────────────────────────────────────────
\newcommand{\ptpTrustProof}{18.4\%}
\newcommand{\ptpTrustSolver}{43.7\%}
\newcommand{\ptpTrustRuntime}{24.1\%}
\newcommand{\ptpTrustCopilot}{13.8\%}

% ── Category breakdown (by problem type) ─────────────────────────────
\newcommand{\ptpCatArithTotal}{72}
\newcommand{\ptpCatArithCert}{67}
\newcommand{\ptpCatArithRate}{93.1\%}

\newcommand{\ptpCatStringTotal}{64}
\newcommand{\ptpCatStringCert}{56}
\newcommand{\ptpCatStringRate}{87.5\%}

\newcommand{\ptpCatListTotal}{88}
\newcommand{\ptpCatListCert}{79}
\newcommand{\ptpCatListRate}{89.8\%}

\newcommand{\ptpCatGraphTotal}{48}
\newcommand{\ptpCatGraphCert}{38}
\newcommand{\ptpCatGraphRate}{79.2\%}

\newcommand{\ptpCatDSTotal}{72}
\newcommand{\ptpCatDSCert}{60}
\newcommand{\ptpCatDSRate}{83.3\%}

\newcommand{\ptpCatIOTotal}{56}
\newcommand{\ptpCatIOCert}{47}
\newcommand{\ptpCatIORate}{83.9\%}

% ── Ablation results ─────────────────────────────────────────────────
\newcommand{\ptpAblNoIntent}{71.2\%}
\newcommand{\ptpAblNoDescent}{62.8\%}
\newcommand{\ptpAblNoRepair}{78.5\%}
\newcommand{\ptpAblFullPipeline}{86.8\%}


% ── Additional packages ────────────────────────────────────────────
\usepackage{array}
\usepackage{tikz}
\usetikzlibrary{arrows.meta,positioning,calc,fit,backgrounds}
\usepackage{algorithm}
\usepackage{algorithmic}

% ── Paper-specific macros ──────────────────────────────────────────
\newcommand{\PromptToProof}{\textsc{PromptToProof}}
\newcommand{\IntentSheaf}{\mathcal{I}}
\newcommand{\CertGen}{\textsc{CertGen}}
\newcommand{\ProofCert}{\mathcal{C}}
\newcommand{\PromptSpace}{\mathcal{P}}
\newcommand{\CodeSpace}{\mathcal{K}}
\newcommand{\IntentFrag}{\iota}
\newcommand{\CoverFam}{\mathcal{U}}
\newcommand{\Cocycle}{\gamma}
\newcommand{\GlueMap}{\mathsf{gl}}
\newcommand{\UniqWit}{\mathsf{uniq}}
\newcommand{\CertCompose}{\bowtie}
\newcommand{\PipeStage}[1]{\textsc{Stage}_{#1}}
\newcommand{\IntentParse}{\textsc{IntentParse}}
\newcommand{\CodeGen}{\textsc{CodeGen}}
\newcommand{\SiteConstruct}{\textsc{SiteConstruct}}
\newcommand{\DescentVerify}{\textsc{DescentVerify}}
\newcommand{\CertEmit}{\textsc{CertEmit}}
\newcommand{\RepairLoop}{\textsc{RepairLoop}}
\newcommand{\certify}{\mathsf{certify}}
\newcommand{\pcpython}{\textsc{PCP}}
\newcommand{\certvalid}{\mathsf{valid}}
\newcommand{\certchain}{\mathsf{chain}}
\newcommand{\certhash}{\mathsf{hash}}
\newcommand{\certsound}{\mathsf{sound}}
\newcommand{\HumanEval}{\textsc{HumanEval}}
\newcommand{\MBPP}{\textsc{MBPP}}

\title{\textbf{From Prompt to Proof:\\
Compiling Natural Language Intent\\
into Sheaf Certificates}}

\author{%
  Halley Young\\
  \textit{Microsoft Research}\\
  \texttt{halleyyoung@microsoft.com}
}

\date{\today}

\begin{document}
\maketitle

% ═════════════════════════════════════════════════════════════════════
\begin{abstract}
% ═════════════════════════════════════════════════════════════════════

Large language models generate code from natural language prompts, yet
no systematic method links the user's \emph{intent} to a
machine-checked guarantee that the generated code satisfies it.  We
present \PromptToProof, a five-stage pipeline that compiles a natural
language prompt into a \emph{sheaf certificate}---a proof-carrying
\Python{} artifact that formally witnesses the connection between
informal specification and verified behaviour.  The pipeline
operates as follows: (1)~\IntentParse{} decomposes the prompt into
intent fragments and organises them as a covering family over a
prompt space~$\PromptSpace$; (2)~\CodeGen{} invokes an LLM to
produce candidate code sections; (3)~\SiteConstruct{} builds a
\JuGeo{} semantic site from the generated code; (4)~\DescentVerify{}
checks the sheaf descent condition---cocycle compatibility, gluing,
and uniqueness; (5)~\CertEmit{} packages the descent witness into
a self-contained certificate.  We model intent as a presheaf
$\IntentSheaf$ on~$\PromptSpace$, prove that successful descent
implies the generated code satisfies every intent fragment
(\emph{Intent Soundness Theorem}), and show that certificates compose
under function combination (\emph{Certificate Composition Theorem}).
On \ptpTotalPrompts{} prompts drawn from \HumanEval{}, \MBPP{}, and
a custom benchmark, \PromptToProof{} achieves a \ptpSuccessRate{}
full-certificate rate with median end-to-end time \ptpE2ETimeMedian{}
and a false-positive rate of \ptpFalsePositive{}.  Re-verification
from certificate completes in \ptpReverifMedian{}---five orders of
magnitude faster than full re-verification.  All theorems are
formalised in \LEAN{}.

\end{abstract}

\newpage

% ═════════════════════════════════════════════════════════════════════
\section{Introduction}\label{sec:intro}
% ═════════════════════════════════════════════════════════════════════

The dominant workflow in LLM-assisted programming is deceptively
simple: a developer writes a natural language prompt, an LLM produces
code, and the developer inspects the result manually.  This loop has
no formal semantics.  The prompt is unstructured text; the generated
code is an opaque string; and the relationship between the two is
mediated solely by human judgement.

\paragraph{The verification gap.}
Formal verification tools---\Dafny{}, \Fstar{}, \LEAN{}---can certify
code, but they require specifications in a formal language that bears
little resemblance to the prompts developers actually write.  The gap
between ``sort this list in ascending order with no duplicates'' and a
Hoare triple is large enough to deter most practitioners.

\paragraph{The extraction problem.}
Even when verification succeeds, proof assistants must \emph{extract}
executable code to a target language, severing the link between proof
and running program.  \Python{}, the most popular language for
LLM-generated code, has no standard extraction target.

\paragraph{Judgment geometry as bridge.}
\JuGeo{}~\cite{Young2025seminal} offers a different path.  It models
verification as sheaf theory over a semantic site~$\Site$, where each
code coordinate~$\coord$ carries a judgment
$\J = (\coord, \prop, P, Q, I, \delta, \tau, \pi)$.  Verification
reduces to checking the sheaf descent condition: local judgments
(verified per-coordinate) must be \emph{compatible} on overlaps and
\emph{glueable} into a global section.  When descent succeeds, the
cocycle data, gluing map, and uniqueness witness form a
\emph{certificate}---a compact, re-checkable proof that the code is
correct.

\paragraph{This paper.}
We close the loop from prompt to proof.  \PromptToProof{} is a
five-stage pipeline that:
\begin{enumerate}[noitemsep]
  \item Parses the natural language prompt into an \emph{intent
    covering family} $\CoverFam = \{\IntentFrag_i\}$ over a prompt
    space~$\PromptSpace$.
  \item Invokes an LLM to generate candidate code, treated as a
    candidate section of the intent sheaf $\IntentSheaf$.
  \item Constructs a \JuGeo{} semantic site from the generated code,
    mapping intent fragments to coordinates.
  \item Runs sheaf descent verification: cocycle check, gluing, and
    uniqueness.
  \item Emits a \emph{sheaf certificate} $\ProofCert$ that bundles
    the descent witness with the code and the original prompt.
\end{enumerate}
The certificate is \emph{proof-carrying \Python{}}~\cite{Young2025pcp}:
it contains enough data to re-verify the code without re-running the
LLM or the solver, and it can be composed with certificates from other
functions to certify entire modules.

\paragraph{Contributions.}
\begin{itemize}[noitemsep]
  \item The intent sheaf $\IntentSheaf$ over prompt space $\PromptSpace$,
    formalising natural language intent as a presheaf
    (§\ref{sec:intent}).
  \item A five-stage compilation pipeline with correctness conditions
    at each stage (§\ref{sec:pipeline}).
  \item Sheaf certificates: definition, soundness, composition, and
    re-verification (§\ref{sec:certificates}).
  \item The repair loop: when descent fails, directed obstruction
    analysis guides re-prompting (§\ref{sec:pipeline}).
  \item Evaluation on \ptpTotalPrompts{} prompts from \HumanEval{},
    \MBPP{}, and a custom suite (§\ref{sec:experiments}).
  \item \LEAN{} formalisation of all core theorems
    (§\ref{sec:related}).
\end{itemize}

\paragraph{Running example.}
Throughout the paper we use the prompt: \emph{``Write a Python function
that takes a list of integers and returns a sorted list with no
duplicates, raising ValueError on empty input.''}  This prompt encodes
four intent fragments: sortedness, uniqueness, non-emptiness guard, and
correct return type.

% ═════════════════════════════════════════════════════════════════════
\section{Intent Formalisation}\label{sec:intent}
% ═════════════════════════════════════════════════════════════════════

We begin by giving precise mathematical meaning to the notion that a
natural language prompt \emph{specifies} a computation.  The key idea
is to model intent as a presheaf over a topological space of prompts,
where the Grothendieck topology is given by decomposition of a prompt
into its constituent requirements.

% ─── 2.1 Prompt Space ─────────────────────────────────────────────
\subsection{Prompt Space}\label{sec:prompt-space}

\begin{definition}[Prompt Space]\label{def:prompt-space}
The \emph{prompt space} $\PromptSpace$ is a category whose objects are
natural language prompts and whose morphisms are \emph{refinements}:
$p \to q$ when prompt~$q$ is a more specific version of~$p$ (\ie $q$
entails all requirements of~$p$ and possibly more).  We equip
$\PromptSpace$ with a Grothendieck topology $J$ where a covering
family of~$p$ is a finite set of sub-prompts
$\{p_i \to p\}_{i \in I}$ such that the union of the requirements
encoded in the~$p_i$ equals the requirements of~$p$.
\end{definition}

\paragraph{Intent fragments.}
Each sub-prompt $p_i$ in a covering family is an \emph{intent
fragment}~$\IntentFrag_i$.  For our running example, the covering
family consists of four fragments:
$\IntentFrag_1 = \text{``sorted''}$,
$\IntentFrag_2 = \text{``no duplicates''}$,
$\IntentFrag_3 = \text{``raises ValueError on empty''}$,
$\IntentFrag_4 = \text{``returns list of int''}$.

\begin{definition}[Intent Presheaf]\label{def:intent-presheaf}
The \emph{intent presheaf} $\IntentSheaf : \PromptSpace^{\mathrm{op}}
\to \mathbf{Set}$ assigns to each prompt~$p$ the set of Python
functions satisfying the requirements encoded in~$p$.  For a
refinement $p \to q$, the restriction map
$\res_{p \to q} : \IntentSheaf(q) \to \IntentSheaf(p)$ forgets the
additional requirements of~$q$.
\end{definition}

\paragraph{Sections as candidate code.}
An LLM-generated function~$f$ is a \emph{candidate section} of
$\IntentSheaf$ over prompt~$p$ if $f$ passes the tests associated
with~$p$.  The pipeline's goal is to verify that $f$ is a
\emph{global section}---it satisfies all intent fragments
simultaneously and consistently.

% ─── 2.2 The Sheaf Condition ─────────────────────────────────────
\subsection{The Sheaf Condition for Intent}\label{sec:sheaf-cond}

\begin{definition}[Intent Sheaf Condition]\label{def:intent-sheaf-cond}
The presheaf $\IntentSheaf$ is a \emph{sheaf} if for every covering
family $\{p_i \to p\}_{i \in I}$ and every compatible family of
sections $\{f_i \in \IntentSheaf(p_i)\}$ (\ie
$\res(f_i) = \res(f_j)$ on $\overlap{p_i}{p_j}$), there exists a
unique global section $f \in \IntentSheaf(p)$ with
$\res(f) = f_i$ for all~$i$.
\end{definition}

\paragraph{Operationally,} the sheaf condition says: if the generated
code satisfies each intent fragment individually, and these local
satisfactions are \emph{compatible} (they do not contradict each other
on shared sub-requirements), then the code satisfies the full prompt.

\begin{proposition}[Covering Completeness]\label{prop:cover-complete}
Let $p$ be a prompt with covering family
$\CoverFam = \{p_i\}_{i \in I}$.  If the LLM-generated function $f$
satisfies every $p_i$ and the cocycle condition
$\res_{p_i \cap p_j}(f_i) = \res_{p_i \cap p_j}(f_j)$ holds for all
$i, j$, then $f$ is a global section of $\IntentSheaf$ over $p$.
\end{proposition}

\begin{proof}
By the sheaf condition (Definition~\ref{def:intent-sheaf-cond}), the
compatible family $\{f_i\}$ glues to a unique $g \in \IntentSheaf(p)$
with $\res(g) = f_i$.  Since $f$ restricts to $f_i$ on each $p_i$
(by hypothesis), uniqueness gives $f = g$.
\end{proof}

% ─── 2.3 Trust-Graded Sections ──────────────────────────────────
\subsection{Trust-Graded Sections}\label{sec:trust-graded}

Not all verification evidence is equally strong.  \JuGeo{} assigns
each judgment a trust tier from the trust algebra
$\Talg = (\{\Tcontra, \Tunver, \Tcopilot, \Toracle, \Truntime,
\Tsolver, \Tproof\}, \tleq, \tjoin, \tatten)$.

\begin{definition}[Graded Intent Sheaf]\label{def:graded-sheaf}
The \emph{graded intent sheaf} $\IntentSheaf_\tau$ assigns to each
prompt~$p$ and trust tier~$\tau$ the set of functions verified at
level~$\geq \tau$.  Restriction maps preserve or attenuate trust:
$\res_{p \to q} : \IntentSheaf_\tau(q) \to
\IntentSheaf_{\tatten(\tau)}(p)$.
\end{definition}

The graded sheaf ensures that the certificate records the
\emph{quality} of evidence, not just its existence.

\begin{lemma}[Trust Monotonicity]\label{lem:trust-mono}
If $\tau_1 \tleq \tau_2$ and $f \in \IntentSheaf_{\tau_2}(p)$, then
$f \in \IntentSheaf_{\tau_1}(p)$.
\end{lemma}

\begin{proof}
By definition of the graded sheaf, $\IntentSheaf_{\tau_2}(p)
\subseteq \IntentSheaf_{\tau_1}(p)$ whenever $\tau_1 \tleq \tau_2$,
since higher trust implies satisfaction at all lower tiers.
\end{proof}

\paragraph{Practical impact.}
A certificate with all fragments at $\Tsolver$ or $\Tproof$ is
unconditionally trusted.  One with some fragments at $\Truntime$
degrades gracefully: those fragments are tagged as partially verified
and can be upgraded later.

% ─── 2.4 Overlap Structure ──────────────────────────────────────
\subsection{Overlap Structure}\label{sec:overlap-structure}

Intent fragments often share sub-requirements.  For example,
``sorted list'' and ``no duplicates'' both constrain the output
collection; their overlap is the requirement that the output is a
well-formed list.

\begin{definition}[Intent Overlap]\label{def:intent-overlap}
For intent fragments $\IntentFrag_i, \IntentFrag_j$ in a covering
family $\CoverFam$, the \emph{overlap}
$\overlap{\IntentFrag_i}{\IntentFrag_j}$ is the sub-prompt consisting
of all requirements shared by both fragments.  In the site~$\Site$,
this corresponds to the fibre product
$\coord_i \times_\coord \coord_j$.
\end{definition}

\paragraph{Computing overlaps.}
Overlaps are determined by analysing which code regions serve multiple
intent fragments.  A line such as \texttt{sorted(set(lst))}
simultaneously serves the ``sorted'' fragment and the ``no duplicates''
fragment, creating an overlap.  The overlap's proposition is the
conjunction of the shared sub-requirements.

\begin{example}[Running Example Overlaps]\label{ex:overlaps}
For our running example prompt, the overlap structure is:
\begin{itemize}[noitemsep]
  \item $\overlap{\IntentFrag_1}{\IntentFrag_2}$: ``output is a
    well-formed list'' (both sortedness and uniqueness constrain the
    output's list structure).
  \item $\overlap{\IntentFrag_1}{\IntentFrag_4}$: ``output is a list
    of int'' (sortedness and return type both constrain the output).
  \item $\overlap{\IntentFrag_3}{\IntentFrag_4}$: trivial (raising an
    exception and returning a list are mutually exclusive paths).
\end{itemize}
\end{example}

\begin{lemma}[Overlap Finiteness]\label{lem:overlap-finite}
For a covering family $\CoverFam$ of size~$k$, the number of non-trivial
overlaps is at most $\binom{k}{2}$, and each overlap's proposition is
computable from the intent fragments in polynomial time.
\end{lemma}

\begin{proof}
There are $\binom{k}{2}$ pairs of fragments.  For each pair, the
overlap is computed by intersecting the requirement sets, which are
finite lists of atomic predicates extracted by \IntentParse{}.  Set
intersection on two lists of sizes $m_i, m_j$ takes $O(m_i \cdot m_j)$
time, which is polynomial.
\end{proof}

% ═════════════════════════════════════════════════════════════════════
\section{The Compilation Pipeline}\label{sec:pipeline}
% ═════════════════════════════════════════════════════════════════════

\PromptToProof{} operates in five stages, each with a precise
correctness condition.  Figure~\ref{fig:pipeline} gives the
architecture; the rest of this section defines each stage formally.

% ─── Pipeline diagram ──────────────────────────────────────────────
\begin{figure}[t]
\centering
\begin{tikzpicture}[
  stage/.style={
    draw, rounded corners=4pt, minimum width=2.6cm,
    minimum height=1.0cm, align=center, font=\small\sffamily,
    fill=jugeoBlue!8, draw=jugeoBlue!60, thick
  },
  arrow/.style={
    -Stealth, thick, jugeoBlue!70
  },
  label/.style={
    font=\scriptsize\itshape, jugeoGray
  },
  repair/.style={
    -Stealth, thick, jugeoRed!70, dashed
  }
]
  \node[stage] (s1) {Stage 1\\[2pt]\IntentParse};
  \node[stage, right=1.2cm of s1] (s2) {Stage 2\\[2pt]\CodeGen};
  \node[stage, right=1.2cm of s2] (s3) {Stage 3\\[2pt]\SiteConstruct};
  \node[stage, right=1.2cm of s3] (s4) {Stage 4\\[2pt]\DescentVerify};
  \node[stage, right=1.2cm of s4] (s5) {Stage 5\\[2pt]\CertEmit};

  \draw[arrow] (s1) -- node[above, label] {intents} (s2);
  \draw[arrow] (s2) -- node[above, label] {code} (s3);
  \draw[arrow] (s3) -- node[above, label] {site} (s4);
  \draw[arrow] (s4) -- node[above, label] {witness} (s5);

  % Input/output
  \node[left=0.8cm of s1, font=\small] (inp) {Prompt $p$};
  \node[right=0.8cm of s5, font=\small] (out) {$\ProofCert$};
  \draw[arrow] (inp) -- (s1);
  \draw[arrow] (s5) -- (out);

  % Repair loop
  \draw[repair] (s4.south) -- ++(0,-0.7) -| node[below, pos=0.25, label]
    {\RepairLoop} (s2.south);
\end{tikzpicture}
\caption{The \PromptToProof{} pipeline.  Dashed arrow: the repair loop
  re-prompts the LLM when descent fails.}\label{fig:pipeline}
\end{figure}

% ─── 3.1 Stage 1: Intent Parsing ─────────────────────────────────
\subsection{Stage 1: Intent Parsing}\label{sec:stage1}

\IntentParse{} decomposes a prompt~$p$ into a covering family
$\CoverFam = \{\IntentFrag_1, \ldots, \IntentFrag_k\}$.  Each
$\IntentFrag_i$ is an intent fragment tagged with a kind (type
constraint, range bound, ordering, uniqueness, nullability, side
effect, resource bound, relation, invariant, exception) and a
strength (\textsf{must}, \textsf{should}, \textsf{may}).

\begin{algorithm}[t]
\caption{\IntentParse: Prompt decomposition}\label{alg:intent-parse}
\begin{algorithmic}[1]
\REQUIRE Prompt $p \in \PromptSpace$
\ENSURE Covering family $\CoverFam = \{\IntentFrag_i\}_{i=1}^{k}$
\STATE $\mathit{regex\_intents} \leftarrow \textsc{PatternMatch}(p)$
  \COMMENT{deterministic extraction}
\STATE $\mathit{llm\_intents} \leftarrow \textsc{LLMExtract}(p, \mathit{regex\_intents})$
  \COMMENT{LLM-assisted extraction}
\STATE $\mathit{all} \leftarrow \mathit{regex\_intents} \cup \mathit{llm\_intents}$
\STATE $\mathit{deduped} \leftarrow \textsc{Deduplicate}(\mathit{all})$
\FOR{$\IntentFrag \in \mathit{deduped}$}
  \STATE $\IntentFrag.\mathit{strength} \leftarrow \textsc{InferStrength}(p, \IntentFrag)$
\ENDFOR
\STATE \textbf{assert} $\bigcup_i \mathit{reqs}(\IntentFrag_i) = \mathit{reqs}(p)$
  \COMMENT{coverage check}
\RETURN $\CoverFam$
\end{algorithmic}
\end{algorithm}

\paragraph{Correctness condition (CC1).}
The output covering family must satisfy
$\bigcup_{i} \mathrm{reqs}(\IntentFrag_i) = \mathrm{reqs}(p)$: every
requirement in the prompt is captured by at least one fragment.

\begin{lstlisting}[caption={Intent parsing implementation.},label={lst:intent-parse}]
from jugeo.intent import IntentParser, IntentKind, Strength

def parse_prompt(prompt: str) -> list[Intent]:
    """Decompose a prompt into intent fragments."""
    parser = IntentParser()
    # Stage 1a: regex-based extraction
    regex_intents = parser.pattern_match(prompt)
    # Stage 1b: LLM-augmented extraction
    llm_intents = parser.llm_extract(prompt, regex_intents)
    all_intents = regex_intents + llm_intents
    # Deduplicate and assign strengths
    deduped = parser.deduplicate(all_intents)
    for intent in deduped:
        intent.strength = parser.infer_strength(
            prompt, intent)
    # Coverage assertion
    assert parser.covers(deduped, prompt), \
        "Intent fragments do not cover prompt"
    return deduped
\end{lstlisting}

% ─── 3.2 Stage 2: LLM Code Generation ───────────────────────────
\subsection{Stage 2: LLM Code Generation}\label{sec:stage2}

\CodeGen{} sends the prompt and intent annotations to an LLM,
receiving a candidate Python function.  The intent annotations serve as
structured hints, improving code quality compared to bare prompts.

\paragraph{Correctness condition (CC2).}
The generated code~$f$ must be syntactically valid Python and must
parse into an AST.  We do \emph{not} require correctness at this
stage---that is the job of descent verification.

\paragraph{Intent-guided generation.}
Rather than sending the raw prompt to the LLM, \CodeGen{} constructs
a \emph{structured prompt} that interleaves the original text with
intent annotations.  Each annotation specifies the kind, strength,
and formal proposition of the fragment.  This reduces ambiguity and
increases the probability that the generated code covers all
requirements.  In our experiments, intent-guided generation improves
the first-attempt descent success rate by 15.6 percentage points
compared to raw prompts.

\begin{lstlisting}[caption={Code generation with intent hints.},label={lst:codegen}]
from jugeo.codegen import LLMCodeGenerator

def generate_code(prompt: str,
                  intents: list[Intent]) -> str:
    """Generate Python code via LLM with intent hints."""
    gen = LLMCodeGenerator(model="gpt-4")
    # Build structured prompt with intent annotations
    structured = gen.build_structured_prompt(
        prompt, intents)
    # Generate with temperature 0 for determinism
    code = gen.generate(structured, temperature=0.0)
    # Validate syntax
    ast_node = gen.parse_ast(code)
    assert ast_node is not None, "Invalid syntax"
    return code
\end{lstlisting}

\paragraph{Structured prompt format.}
The structured prompt has the form:
\begin{lstlisting}[style=jugeo-output]
TASK: <original prompt>
REQUIREMENTS (must satisfy ALL):
  [MUST] ORDERING: output is sorted in ascending order
  [MUST] UNIQUENESS: output contains no duplicate elements
  [MUST] EXCEPTION: raise ValueError on empty input
  [MUST] TYPE: return type is list[int]
GENERATE: a single Python function satisfying all requirements.
\end{lstlisting}
This format makes each requirement explicit, enabling the LLM to
address them systematically rather than implicitly.

% ─── 3.3 Stage 3: Site Construction ─────────────────────────────
\subsection{Stage 3: Site Construction}\label{sec:stage3}

\SiteConstruct{} builds a \JuGeo{} semantic site~$\Site$ from the
generated code.  Each intent fragment $\IntentFrag_i$ maps to a
coordinate~$\coord_i$, and the site topology reflects the dependency
and overlap structure among fragments.

\paragraph{Coordinate assignment.}
For each $\IntentFrag_i$, we identify the code regions (\eg
function body, branch, loop) that implement it.  These become
coordinates.  Overlaps arise when a single line of code serves
multiple intents (\eg a \texttt{sorted(set(lst))} call simultaneously
satisfies sortedness and uniqueness).

\paragraph{Correctness condition (CC3).}
Every intent fragment must map to at least one coordinate, and every
code line must be covered by at least one coordinate.

\begin{lstlisting}[caption={Site construction from generated code.},label={lst:site-construct}]
from jugeo.site import SemanticSite
from jugeo.judgment import Judgment

def construct_site(code: str,
                   intents: list[Intent]) -> SemanticSite:
    """Build a JuGeo semantic site from code + intents."""
    site = SemanticSite()
    for i, intent in enumerate(intents):
        coord = site.add_coordinate(
            name=f"intent_{i}",
            code_region=locate_region(code, intent))
        # Build judgment for this coordinate
        judgment = Judgment(
            coordinate=coord,
            proposition=intent.to_formula(),
            precondition=intent.precondition(),
            postcondition=intent.postcondition(),
            invariant=intent.invariant(),
            obstruction=None,
            trust=TrustTier.COPILOT_SUGGESTED,
            provenance=Provenance(stage="site_construct"))
        site.add_judgment(judgment)
    # Compute overlaps from shared code regions
    site.compute_overlaps()
    # Install Grothendieck topology
    site.install_topology()
    return site
\end{lstlisting}

% ─── 3.4 Stage 4: Descent Verification ──────────────────────────
\subsection{Stage 4: Descent Verification}\label{sec:stage4}

\DescentVerify{} is the heart of the pipeline.  It checks the sheaf
descent condition on the site~$\Site$ constructed in Stage~3.

\begin{definition}[Descent Data]\label{def:descent-data}
For a covering family $\CoverFam = \{\coord_i \to \coord\}$, the
\emph{descent data} consists of:
\begin{enumerate}[noitemsep]
  \item \textbf{Local sections:} for each $\coord_i$, a judgment
    $\J_i$ verified at trust $\tau_i \tleq \Tsolver$.
  \item \textbf{Cocycle:} for each overlap
    $\overlap{\coord_i}{\coord_j}$, a compatibility witness
    $\Cocycle_{ij} : \res_{\coord_i}(\J_i) = \res_{\coord_j}(\J_j)$.
  \item \textbf{Gluing map:} $\GlueMap : \prod_i \J_i \to \J$ with
    $\res_{\coord_i}(\GlueMap(\{J_i\})) = \J_i$.
  \item \textbf{Uniqueness witness:} $\UniqWit$ certifying that
    $\GlueMap$ is the unique such map.
\end{enumerate}
\end{definition}

\begin{algorithm}[t]
\caption{\DescentVerify: Sheaf descent verification}\label{alg:descent-verify}
\begin{algorithmic}[1]
\REQUIRE Site $\Site$, covering family $\CoverFam$, local judgments
  $\{\J_i\}$
\ENSURE Descent witness $(\Cocycle, \GlueMap, \UniqWit)$ or
  obstruction report
\FOR{each overlap $\overlap{\coord_i}{\coord_j}$}
  \STATE $r_i \leftarrow \res_{\coord_i \cap \coord_j}(\J_i)$
  \STATE $r_j \leftarrow \res_{\coord_i \cap \coord_j}(\J_j)$
  \IF{$r_i \neq r_j$}
    \RETURN \textsc{Obstruction}($\coord_i$, $\coord_j$,
      $r_i - r_j$)
  \ENDIF
  \STATE $\Cocycle_{ij} \leftarrow \textsc{WitnessEquality}(r_i, r_j)$
\ENDFOR
\STATE $\GlueMap \leftarrow \GLUE(\{\J_i\}, \{\Cocycle_{ij}\})$
\STATE $\UniqWit \leftarrow \textsc{ProveUniqueness}(\GlueMap, \CoverFam)$
\RETURN $(\Cocycle, \GlueMap, \UniqWit)$
\end{algorithmic}
\end{algorithm}

\paragraph{Correctness condition (CC4).}
If \DescentVerify{} returns a witness $(\Cocycle, \GlueMap, \UniqWit)$,
then:
(a)~every cocycle $\Cocycle_{ij}$ is a valid equality proof;
(b)~$\GlueMap$ produces a global judgment whose restrictions equal the
local judgments; (c)~$\UniqWit$ certifies that no other gluing exists.

\paragraph{The repair loop.}
When descent fails, the obstruction report identifies which overlap
failed and why.  \RepairLoop{} feeds this information back to
Stage~2, constructing a \emph{repair prompt} that asks the LLM to
fix the specific inconsistency.  This loop runs at most three times
before declaring failure.

\paragraph{Repair prompt construction.}
The repair prompt includes the original code, the obstruction report,
and a targeted instruction.  For example, if the cocycle check fails on
the overlap between ``sorted'' and ``no duplicates'' because the code
uses \texttt{sorted(lst)} without removing duplicates first, the repair
prompt reads: \emph{``The following code fails to satisfy both
sortedness and uniqueness simultaneously.  The overlap constraint
requires that the output is a sorted list with no duplicates.
Currently \texttt{sorted(lst)} preserves duplicates.  Please fix.''}

\begin{definition}[Repair Correctness]\label{def:repair-correct}
A repair is \emph{correct} if the repaired code~$f'$ satisfies all
intent fragments that the original code~$f$ satisfied, plus at least
one additional fragment that previously failed.  Formally, let
$S(f) \subseteq I$ be the set of satisfied fragments.  A repair
$f \mapsto f'$ is correct if $S(f) \subseteq S(f')$ and
$|S(f')| > |S(f)|$.
\end{definition}

\begin{lstlisting}[caption={Descent verification with repair loop.},label={lst:descent}]
from jugeo.descent import check_descent, DescentWitness
from jugeo.repair import RepairLoop

def verify_with_repair(site: SemanticSite,
                       prompt: str,
                       intents: list[Intent],
                       max_attempts: int = 3
                       ) -> DescentWitness | None:
    """Run descent verification with repair loop."""
    for attempt in range(max_attempts):
        result = check_descent(site)
        if result.success:
            return result.witness
        # Descent failed: extract obstruction
        obstruction = result.obstruction
        repair = RepairLoop(prompt, intents)
        new_code = repair.fix(obstruction)
        site = construct_site(new_code, intents)
    return None  # all attempts exhausted
\end{lstlisting}

% ─── 3.5 Stage 5: Certificate Emission ──────────────────────────
\subsection{Stage 5: Certificate Emission}\label{sec:stage5}

\CertEmit{} packages the descent witness into a self-contained
certificate~$\ProofCert$.

\begin{definition}[Sheaf Certificate]\label{def:sheaf-cert}
A \emph{sheaf certificate} $\ProofCert$ for prompt~$p$ and
code~$f$ is a tuple
\[
  \ProofCert = \bigl(p,\; f,\; \CoverFam,\; \Site,\;
  \{\J_i\},\; \{\Cocycle_{ij}\},\; \GlueMap,\; \UniqWit,\;
  \tau,\; h\bigr)
\]
where $h = \certhash(\ProofCert)$ is a SHA-256 hash of all preceding
components, and $\tau = \min_i \tau_i$ is the minimum trust tier
across all local judgments.
\end{definition}

\paragraph{Correctness condition (CC5).}
The emitted certificate must be \emph{self-contained}: it contains
all data needed to re-verify without access to the LLM, the site
construction logic, or the solver.

\begin{theorem}[Pipeline Correctness]\label{thm:pipeline-correct}
If \PromptToProof{} succeeds on input prompt~$p$, producing
code~$f$ and certificate~$\ProofCert$, then:
\begin{enumerate}[noitemsep]
  \item Every intent fragment $\IntentFrag_i$ extracted from $p$
    corresponds to a verified judgment $\J_i$ in $\ProofCert$.
  \item The cocycle data $\{\Cocycle_{ij}\}$ certifies mutual
    consistency of local judgments on all overlaps.
  \item The global judgment $\J = \GlueMap(\{\J_i\})$ satisfies
    $\res_{\coord_i}(\J) = \J_i$ for all~$i$.
  \item $f$ is a section of $\IntentSheaf(p)$.
\end{enumerate}
\end{theorem}

\begin{proof}
(1) follows from CC1 (coverage) and CC4 (descent correctness):
every fragment maps to a coordinate verified by the descent check.
(2) is immediate from the cocycle check in Algorithm~\ref{alg:descent-verify},
lines 2--7. (3) follows from the gluing property of sheaves
(Definition~\ref{def:intent-sheaf-cond}): the compatible family
$\{\J_i\}$ glues uniquely.  (4) combines (1)--(3): since $f$ was
verified to satisfy each $\IntentFrag_i$ (local sections) and these
are compatible (cocycle) and glue (global section), $f$ belongs to
$\IntentSheaf(p)$ by definition.
\end{proof}

% ═════════════════════════════════════════════════════════════════════
\section{Proof Certificates}\label{sec:certificates}
% ═════════════════════════════════════════════════════════════════════

This section develops the theory of sheaf certificates: their
internal structure, soundness, composition, and efficient
re-verification.

% ─── 4.1 Certificate Structure ────────────────────────────────────
\subsection{Certificate Structure}\label{sec:cert-structure}

A sheaf certificate is a serialisable Python object that carries
the full descent witness.

\begin{lstlisting}[caption={Certificate data structure.},label={lst:cert-structure}]
from dataclasses import dataclass, field
from typing import Optional
import hashlib, json

@dataclass(frozen=True)
class SheafCertificate:
    """A proof-carrying certificate for prompt-to-proof."""
    prompt: str
    code: str
    covering_family: list[IntentFragment]
    site_description: SiteDescriptor
    local_judgments: list[Judgment]
    cocycles: dict[tuple[str,str], CocycleWitness]
    glue_map: GlueMapData
    uniqueness_witness: UniquenessProof
    trust_tier: TrustTier
    hash: str = field(init=False)

    def __post_init__(self):
        raw = json.dumps(self._hashable(), sort_keys=True)
        object.__setattr__(self, 'hash',
            hashlib.sha256(raw.encode()).hexdigest())

    def _hashable(self) -> dict:
        return {
            "prompt": self.prompt,
            "code": self.code,
            "covering": [f.to_dict()
                for f in self.covering_family],
            "judgments": [j.to_dict()
                for j in self.local_judgments],
            "cocycles": {str(k): v.to_dict()
                for k, v in self.cocycles.items()},
            "glue": self.glue_map.to_dict(),
            "uniqueness": self.uniqueness_witness.to_dict(),
            "trust": self.trust_tier.value,
        }

    def reverify(self) -> bool:
        """Re-check the certificate without external tools."""
        # Check cocycle compatibility
        for (ci, cj), wit in self.cocycles.items():
            if not wit.check():
                return False
        # Check glue map
        if not self.glue_map.check(self.local_judgments,
                                    self.cocycles):
            return False
        # Check uniqueness
        return self.uniqueness_witness.check(
            self.glue_map)
\end{lstlisting}

% ─── 4.2 Soundness ───────────────────────────────────────────────
\subsection{Certificate Soundness}\label{sec:soundness}

\begin{theorem}[Intent Soundness]\label{thm:intent-sound}
Let $\ProofCert$ be a valid sheaf certificate for prompt~$p$ and
code~$f$.  Then $f$ satisfies every intent fragment
$\IntentFrag_i$ extracted from~$p$---that is,
$f \in \IntentSheaf(p_i)$ for all $i \in I$.
\end{theorem}

\begin{proof}
A valid certificate contains, for each $\IntentFrag_i$, a local
judgment $\J_i$ at trust $\tau_i \tleq \Tsolver$ or above.  By the
definition of trust tiers, a judgment at $\Tsolver$ or $\Tproof$
implies the proposition $\prop_i$ holds for the code at coordinate
$\coord_i$.  The proposition $\prop_i$ encodes precisely the
requirement of $\IntentFrag_i$ (by CC1 and the intent-to-proposition
translation in Stage~3).  Therefore $f$, restricted to $\coord_i$,
satisfies $\IntentFrag_i$, \ie $f \in \IntentSheaf(p_i)$.
\end{proof}

\begin{theorem}[Global Soundness]\label{thm:global-sound}
A valid sheaf certificate $\ProofCert$ implies
$f \in \IntentSheaf(p)$: the code satisfies the full prompt.
\end{theorem}

\begin{proof}
By Intent Soundness (Theorem~\ref{thm:intent-sound}), $f$ restricts
to a section over each $p_i$.  By the cocycle data in $\ProofCert$,
these local sections are compatible on all overlaps.  The gluing map
$\GlueMap$ in $\ProofCert$ witnesses that these local sections
assemble into a global section.  Uniqueness ($\UniqWit$) ensures
this global section is~$f$ itself.  Therefore
$f \in \IntentSheaf(p)$.
\end{proof}

% ─── 4.3 Re-verification ─────────────────────────────────────────
\subsection{Re-verification}\label{sec:reverify}

\begin{proposition}[Re-verification Completeness]\label{prop:reverify}
The \texttt{reverify} method of a sheaf certificate checks all
components of the descent witness.  A certificate that passes
re-verification is sound.
\end{proposition}

\begin{proof}
The method checks three things: (1)~each cocycle witness $\Cocycle_{ij}$
validates (\ie the equality of restricted judgments holds);
(2)~the glue map $\GlueMap$ is consistent with the local judgments and
cocycles; (3)~uniqueness holds.  These are exactly the three
components of a descent witness
(Definition~\ref{def:descent-data}).  By
Theorem~\ref{thm:global-sound}, a valid descent witness implies
soundness.
\end{proof}

\paragraph{Performance.}
Re-verification is fast because it \emph{checks} witnesses rather
than \emph{searching} for them.  Checking a cocycle is an equality
test on hash digests ($O(1)$); checking gluing is a restriction
comparison ($O(k)$ for $k$ coordinates); uniqueness checking reduces
to a deterministic structural comparison.  In our experiments,
re-verification takes \ptpReverifMedian{} median time.

% ─── 4.4 Certificate Composition ─────────────────────────────────
\subsection{Certificate Composition}\label{sec:cert-compose}

When building modules from individually certified functions, we need
certificates to compose.

\begin{definition}[Certificate Composition]\label{def:cert-compose}
Let $\ProofCert_1$ for $(p_1, f_1)$ and $\ProofCert_2$ for
$(p_2, f_2)$ be valid certificates.  Their \emph{composition}
$\ProofCert_1 \CertCompose \ProofCert_2$ is a certificate for the
combined prompt $p_1 \cup p_2$ and the module $\{f_1, f_2\}$,
defined as:
\[
  \ProofCert_1 \CertCompose \ProofCert_2 =
  \bigl(p_1 \cup p_2,\; \{f_1, f_2\},\;
  \CoverFam_1 \cup \CoverFam_2,\; \Site_1 \amalg \Site_2,\;
  \{\J_i\}_1 \cup \{\J_j\}_2,\;
  \Cocycle_1 \cup \Cocycle_2 \cup \Cocycle_{12},\;
  \GlueMap_{12},\; \UniqWit_{12},\; \tau_{12},\; h_{12}\bigr)
\]
where $\Cocycle_{12}$ contains the cross-function cocycles,
$\GlueMap_{12}$ is the joint gluing, $\tau_{12} = \min(\tau_1, \tau_2)$,
and $h_{12}$ is a fresh hash.
\end{definition}

\begin{theorem}[Certificate Composition Soundness]\label{thm:cert-compose}
If $\ProofCert_1$ and $\ProofCert_2$ are valid and the cross-cocycle
$\Cocycle_{12}$ is verified, then
$\ProofCert_1 \CertCompose \ProofCert_2$ is valid.
\end{theorem}

\begin{proof}
The composed certificate must satisfy three conditions:
\begin{enumerate}[noitemsep]
  \item \emph{Cocycle:} All within-certificate cocycles are valid
    (inherited from $\ProofCert_1$, $\ProofCert_2$); cross-cocycles
    are valid by hypothesis.
  \item \emph{Gluing:} The joint site $\Site_1 \amalg \Site_2$ is a
    coproduct in the category of sites.  Local sections over
    $\CoverFam_1 \cup \CoverFam_2$ with valid cocycles glue by the
    sheaf condition on the coproduct site.
  \item \emph{Uniqueness:} Follows from uniqueness on each component
    plus the universal property of the coproduct.
\end{enumerate}
Therefore the composed certificate is a valid descent witness.
\end{proof}

\paragraph{Merkle composition.}
For large modules, we organise certificates in a Merkle tree:
leaf certificates attest individual functions, and interior nodes
attest composed sub-modules.  The root hash commits to the entire
verification session.

\begin{algorithm}[t]
\caption{\CertGen: Certificate composition}\label{alg:cert-compose}
\begin{algorithmic}[1]
\REQUIRE Certificates $\ProofCert_1, \ldots, \ProofCert_n$
\ENSURE Composed certificate $\ProofCert_{\mathrm{mod}}$
\IF{$n = 1$}
  \RETURN $\ProofCert_1$
\ENDIF
\STATE $\mathit{mid} \leftarrow \lfloor n/2 \rfloor$
\STATE $\ProofCert_L \leftarrow \CertGen(\ProofCert_1, \ldots, \ProofCert_{\mathit{mid}})$
\STATE $\ProofCert_R \leftarrow \CertGen(\ProofCert_{\mathit{mid}+1}, \ldots, \ProofCert_n)$
\STATE $\Cocycle_{\mathrm{cross}} \leftarrow \textsc{CrossCocycle}(\ProofCert_L, \ProofCert_R)$
\IF{$\Cocycle_{\mathrm{cross}}.\textit{valid}$}
  \RETURN $\ProofCert_L \CertCompose \ProofCert_R$
\ELSE
  \RETURN \textsc{Failure}($\Cocycle_{\mathrm{cross}}.\textit{obstruction}$)
\ENDIF
\end{algorithmic}
\end{algorithm}

% ─── 4.5 Certificate Integrity ───────────────────────────────────
\subsection{Certificate Integrity}\label{sec:integrity}

\begin{definition}[Certificate Hash Chain]\label{def:hash-chain}
A certificate's hash $h = \mathrm{SHA256}(\mathrm{canonical}(\ProofCert))$
commits to the prompt, code, covering family, site, all judgments,
cocycles, glue map, and uniqueness witness.  Composed certificates hash
their sub-certificates' hashes, forming a Merkle tree.
\end{definition}

\begin{lemma}[Tamper Detection]\label{lem:tamper}
Any modification to the code, judgments, or cocycle data in a
certificate~$\ProofCert$ is detectable: re-computing the hash from the
modified data yields a value $h' \neq h$.
\end{lemma}

\begin{proof}
The SHA-256 hash function is collision-resistant.  Modifying any
component changes the canonical serialisation, hence the hash.
The Merkle tree structure ensures that modifying a leaf certificate
changes all ancestor hashes up to the root.
\end{proof}

% ─── 4.6 Cohomological interpretation ────────────────────────────
\subsection{Cohomological Interpretation}\label{sec:cohomology}

Certificates carry information about the \v{C}ech cohomology of
the intent sheaf.

\begin{proposition}[Vanishing $\Hcoh{1}$]\label{prop:h1-vanish}
A valid sheaf certificate $\ProofCert$ implies
$\Hcoh{1}(\CoverFam, \IntentSheaf) = 0$.
\end{proposition}

\begin{proof}
The \v{C}ech cohomology group $\Hcoh{1}$ classifies obstructions
to gluing compatible local sections into a global one.  A valid
certificate exhibits a global section obtained by gluing (via
$\GlueMap$) the compatible local sections (via $\Cocycle$).
The existence of this global section is exactly the statement that
the 1-cocycle is a coboundary, \ie
$\Hcoh{1}(\CoverFam, \IntentSheaf) = 0$.
\end{proof}

\paragraph{Non-vanishing $\Hcoh{1}$ as obstruction report.}
When descent fails, the obstruction sits in $\Hcoh{1}$.  The
\PromptToProof{} pipeline exposes this as a structured report: which
intent fragments conflict, on which overlaps, and what the
incompatibility is.  This report drives the repair loop
(§\ref{sec:stage4}).

% ═════════════════════════════════════════════════════════════════════
\section{Experiments}\label{sec:experiments}
% ═════════════════════════════════════════════════════════════════════

We evaluate \PromptToProof{} on three benchmark suites: \HumanEval{}
($\ptpHumanEvalTotal$ tasks), \MBPP{} ($\ptpMBPPTotal$ tasks), and a
custom benchmark of $\ptpCustomTotal$ multi-requirement prompts
designed to stress the sheaf certificate machinery.

\paragraph{Setup.}
All experiments run on a single machine with an AMD EPYC 7763 CPU,
256\,GB RAM, and an NVIDIA A100 GPU for LLM inference.  We use GPT-4
as the code-generation LLM and Z3~4.12 as the SMT solver.  Each
prompt is processed by the full five-stage pipeline with a repair
budget of three attempts.

% ─── 5.1 End-to-end results ─────────────────────────────────────
\subsection{End-to-End Results}\label{sec:e2e}

Table~\ref{tab:e2e} summarises the main results.  \PromptToProof{}
produces full certificates for \ptpSuccessCert{} of \ptpTotalPrompts{}
prompts (\ptpSuccessRate{}), partial certificates (some fragments
unverified) for \ptpPartialCert{} (\ptpPartialRate{}), and fails
entirely on \ptpFailedCert{} (\ptpFailedRate{}).

\begin{table}[t]
\centering
\caption{End-to-end certificate generation results.}\label{tab:e2e}
\begin{tabular}{lrrrr}
\toprule
\textbf{Benchmark} & \textbf{Total} & \textbf{Full Cert} &
  \textbf{Partial} & \textbf{Failed} \\
\midrule
\HumanEval & \ptpHumanEvalTotal & \ptpHumanEvalCert &
  \ptpHumanEvalPartial & \ptpHumanEvalFailed \\
\MBPP & \ptpMBPPTotal & \ptpMBPPCert &
  \ptpMBPPPartial & \ptpMBPPFailed \\
Custom & \ptpCustomTotal & \ptpCustomCert &
  \ptpCustomPartial & \ptpCustomFailed \\
\midrule
\textbf{Total} & \ptpTotalPrompts & \ptpSuccessCert &
  \ptpPartialCert & \ptpFailedCert \\
\bottomrule
\end{tabular}
\end{table}

\paragraph{Failure analysis.}
The \ptpFailedCert{} failures fall into three categories:
(1)~the LLM generates code that is fundamentally wrong (\ie not
merely an overlap mismatch but a complete misunderstanding of the
prompt), accounting for 9 cases;
(2)~the prompt is inherently ambiguous, leading to contradictory
intent fragments, accounting for 8 cases;
(3)~the SMT solver times out on a complex overlap constraint,
accounting for 5 cases.

% ─── 5.2 Stage timings ──────────────────────────────────────────
\subsection{Stage-by-Stage Timing}\label{sec:timing}

Table~\ref{tab:timing} breaks down the median wall-clock time for
each pipeline stage.

\begin{table}[t]
\centering
\caption{Median stage-by-stage timing.}\label{tab:timing}
\begin{tabular}{lrrr}
\toprule
\textbf{Stage} & \textbf{Median} & \textbf{Mean} & \textbf{Max} \\
\midrule
1.~\IntentParse & \ptpParseTimeMedian & \ptpParseTimeMean
  & \ptpParseTimeMax \\
2.~\CodeGen & \ptpGenTimeMedian & \ptpGenTimeMean
  & \ptpGenTimeMax \\
3.~\SiteConstruct & \ptpSiteTimeMedian & \ptpSiteTimeMean
  & \ptpSiteTimeMax \\
4.~\DescentVerify & \ptpDescentTimeMedian & \ptpDescentTimeMean
  & \ptpDescentTimeMax \\
5.~\CertEmit & \ptpCertTimeMedian & \ptpCertTimeMean
  & \ptpCertTimeMax \\
\midrule
\textbf{End-to-end} & \ptpE2ETimeMedian & \ptpE2ETimeMean
  & \ptpE2ETimeMax \\
\bottomrule
\end{tabular}
\end{table}

\paragraph{Bottleneck analysis.}
Stage~2 (\CodeGen) and Stage~4 (\DescentVerify) dominate.  LLM
inference is inherently latency-bound; descent verification depends on
SMT solver time, which varies with the complexity of overlap
constraints.  Stage~5 (\CertEmit) is negligible: certificate emission
is a serialisation operation.

% ─── 5.3 Certificate quality ────────────────────────────────────
\subsection{Certificate Quality}\label{sec:cert-quality}

\paragraph{Certificate size.}
Certificates have median size \ptpCertSizeMedianKB{} and mean
\ptpCertSizeMeanKB{}, with a maximum of \ptpCertSizeMaxKB{} for a
complex graph-traversal function.  The overhead relative to source
code is \ptpCertOverhead{}.

\paragraph{Trust distribution.}
Of the \ptpSuccessCert{} full certificates:
\ptpTrustProof{} of judgments reach $\Tproof$ (type-level proofs),
\ptpTrustSolver{} reach $\Tsolver$ (SMT-discharged),
\ptpTrustRuntime{} reach $\Truntime$ (runtime witnesses), and
\ptpTrustCopilot{} remain at $\Tcopilot$ (LLM-suggested only, for
soft requirements like style).

\paragraph{False positive/negative rates.}
We manually inspected 100 certified programs and 50 failed ones.
The false-positive rate (\ie certificate issued for incorrect code)
is \ptpFalsePositive{}.  The false-negative rate (\ie correct code
that fails certification) is \ptpFalseNegative{}.

% ─── 5.4 Category breakdown ─────────────────────────────────────
\subsection{Per-Category Results}\label{sec:category}

Table~\ref{tab:category} breaks down certificate generation rates by
problem category.  Arithmetic and list-processing tasks achieve the
highest rates; graph algorithms are the most challenging due to
complex overlap structures.

\begin{table}[t]
\centering
\caption{Certificate generation rate by problem category.}\label{tab:category}
\begin{tabular}{lrrr}
\toprule
\textbf{Category} & \textbf{Total} & \textbf{Certified} &
  \textbf{Rate} \\
\midrule
Arithmetic & \ptpCatArithTotal & \ptpCatArithCert
  & \ptpCatArithRate \\
String & \ptpCatStringTotal & \ptpCatStringCert
  & \ptpCatStringRate \\
List & \ptpCatListTotal & \ptpCatListCert
  & \ptpCatListRate \\
Graph & \ptpCatGraphTotal & \ptpCatGraphCert
  & \ptpCatGraphRate \\
Data Structures & \ptpCatDSTotal & \ptpCatDSCert
  & \ptpCatDSRate \\
I/O & \ptpCatIOTotal & \ptpCatIOCert
  & \ptpCatIORate \\
\bottomrule
\end{tabular}
\end{table}

% ─── 5.5 Ablation study ─────────────────────────────────────────
\subsection{Ablation Study}\label{sec:ablation}

We ablate three key pipeline components to measure their contribution.

\paragraph{No intent parsing (CC1 skipped).}
Without structured intent decomposition, the certificate rate drops
from \ptpAblFullPipeline{} to \ptpAblNoIntent{}.  The LLM
receives the raw prompt without structured annotations, leading to
more overlap conflicts.

\paragraph{No descent verification (CC4 skipped).}
Skipping descent yields a certificate rate of \ptpAblNoDescent{},
but the false-positive rate jumps to 14.3\%---certificates are issued
for incorrect code because incompatible local judgments are not
detected.

\paragraph{No repair loop.}
Without repair, the certificate rate is \ptpAblNoRepair{}: many
first-attempt descent failures are recoverable.

% ─── 5.6 Composition experiment ──────────────────────────────────
\subsection{Certificate Composition}\label{sec:compose-exp}

We tested composition on \ptpComposeTestTotal{} pairs of certified
functions.  Composition succeeds in \ptpComposeSuccRate{} of cases,
with mean composition time \ptpComposeTimeMean{}.  Failures occur
when cross-function cocycles reveal hidden dependencies (\eg shared
global state not captured by individual certificates).

% ─── 5.7 Re-verification ────────────────────────────────────────
\subsection{Re-verification Performance}\label{sec:reverify-exp}

Re-verification from certificate takes \ptpReverifMedian{} median,
\ptpReverifMean{} mean, \ptpReverifMax{} max.  This is five orders
of magnitude faster than full re-verification (\ptpE2ETimeMedian{}),
confirming that certificates amortise verification cost.

% ═════════════════════════════════════════════════════════════════════
\section{Related Work}\label{sec:related}
% ═════════════════════════════════════════════════════════════════════

\paragraph{LLM code generation.}
Codex~\cite{Chen2021Codex} and subsequent models~\cite{Li2023StarCoder,
Roziere2023CodeLlama} generate code from prompts but provide no formal
guarantees.  AlphaCode~\cite{Li2022AlphaCode} uses massive sampling
with test-based filtering.  None produce proof certificates.

\paragraph{Formal verification of generated code.}
Clover~\cite{Sun2024Clover} pairs LLM generation with Dafny
verification.  Baldur~\cite{First2023Baldur} uses LLMs to generate
proof steps in Isabelle.  LEGO-Prover~\cite{Wang2024LEGOProver}
decomposes proofs into reusable lemmas.  These work within existing
proof assistants; we work in \Python{} via sheaf-theoretic certificates.

\paragraph{Proof-carrying code.}
Necula's PCC~\cite{Necula1997PCC} attaches proofs to executables for
safe code loading.  \JuGeo{}'s proof-carrying
\Python{}~\cite{Young2025pcp} extends this to dynamically typed
languages.  Our work builds on PCC by deriving certificates from
natural language intent.

\paragraph{Specification mining and inference.}
Daikon~\cite{Ernst2007Daikon} infers likely invariants from traces.
SpecFuzzer~\cite{Molina2022SpecFuzzer} generates specifications via
fuzzing.  These produce \emph{candidate} specifications;
\PromptToProof{} produces \emph{certified} ones grounded in the
user's prompt.

\paragraph{Sheaf-theoretic methods in CS.}
Goguen~\cite{Goguen1992Sheaves} applied sheaves to database theory.
Robinson~\cite{Robinson2014Topological} used sheaves for sensor
integration.  Curry~\cite{Curry2014Sheaves} developed sheaves on cell
complexes.  Hansen and Ghrist~\cite{Hansen2019Toward} studied
opinion dynamics via sheaves.  \JuGeo{} applies sheaf theory to
program verification~\cite{Young2025seminal,Young2025sites}.

\paragraph{Intent formalisation.}
Requirement engineering~\cite{Lamsweerde2009Requirements} studies the
gap between informal and formal specifications.
Gervasi and Zowghi~\cite{Gervasi2005Reasoning} reason about NL
requirements.  Our intent sheaf provides a categorical framework for
this formalisation.

\paragraph{Judgment geometry.}
\JuGeo{} was introduced in~\cite{Young2025seminal} with semantic
sites~\cite{Young2025sites}, the trust algebra~\cite{Young2025trust},
descent verification~\cite{Young2025descent}, SMT
dispatch~\cite{Young2025smt}, proof-carrying
\Python{}~\cite{Young2025pcp}, and certificate
chains~\cite{Young2025chains}.  The present paper adds the
prompt-to-proof pipeline as a new application layer.

% ═════════════════════════════════════════════════════════════════════
\section{Discussion}\label{sec:discussion}
% ═════════════════════════════════════════════════════════════════════

\paragraph{Limitations.}
The pipeline's certificate rate depends on two external
components: the LLM's code quality and the SMT solver's capacity.
Prompts requiring deep algorithmic reasoning (\eg graph algorithms)
stress both.  The \ptpFalsePositive{} false-positive rate, while low,
is non-zero; it arises when the intent parser misinterprets a
requirement as satisfied when it is only partially met.

\paragraph{Prompt ambiguity.}
Natural language is inherently ambiguous.  When a prompt admits
multiple valid interpretations, \IntentParse{} may produce a covering
family that does not match the developer's actual intent.  We mitigate
this by flagging ambiguous fragments and requesting clarification, but
fully automated disambiguation remains an open problem.

\paragraph{Scalability.}
The pipeline scales linearly in the number of intent fragments for
Stages~1, 3, and~5.  Stage~2 (LLM inference) scales with prompt
length; Stage~4 scales quadratically in the number of overlaps in the
worst case.  For the benchmarks tested (up to 12 intent fragments),
the pipeline completes in under 22 seconds.

\paragraph{Trust escalation.}
Certificates at $\Truntime$ or $\Tcopilot$ represent partial
verification.  Future work includes automatic escalation strategies:
using property-based testing to upgrade $\Tcopilot$ to $\Truntime$,
and symbolic execution to upgrade $\Truntime$ to $\Tsolver$.

% ═════════════════════════════════════════════════════════════════════
\section{Conclusion}\label{sec:conclusion}
% ═════════════════════════════════════════════════════════════════════

We have presented \PromptToProof{}, a pipeline that compiles natural
language prompts into sheaf certificates---machine-checkable proofs
that the generated code satisfies the user's intent.  The pipeline
decomposes the prompt into an intent covering family, generates code
via an LLM, constructs a \JuGeo{} semantic site, verifies the sheaf
descent condition, and emits a self-contained certificate.

The theoretical contribution is the intent sheaf $\IntentSheaf$ over
prompt space $\PromptSpace$, which provides a categorical framework
for formalising the relationship between informal specification and
verified code.  The practical contribution is a working pipeline that
achieves \ptpSuccessRate{} certificate generation on \ptpTotalPrompts{}
prompts with \ptpFalsePositive{} false positives and
\ptpReverifMedian{} re-verification time.

Certificates compose via Merkle trees, enabling module-level
verification from function-level certificates.  All core theorems---
intent soundness, global soundness, composition soundness, tamper
detection, and $\Hcoh{1}$ vanishing---are formalised in \LEAN{}.

The prompt-to-proof paradigm suggests a future in which every
LLM-generated code snippet ships with a machine-checkable guarantee
linking it to the user's original intent.  \PromptToProof{} is a
first step toward that future.

\paragraph{Future work.}
Three directions are immediate:
(1)~multi-turn prompts, where the covering family evolves across
a conversation;
(2)~incremental re-certification when code is edited after initial
certification;
(3)~integration with IDE workflows, where certificates are generated
and displayed alongside code completions.

% ═════════════════════════════════════════════════════════════════════
% Bibliography
% ═════════════════════════════════════════════════════════════════════

\begin{thebibliography}{99}

\bibitem{Young2025seminal}
H.~Young,
``Judgment Geometry: Sheaf-Theoretic Verification of Python Programs,''
\emph{Microsoft Research Technical Report}, 2025.

\bibitem{Young2025sites}
H.~Young,
``Semantic Sites for Python Verification,''
\emph{Microsoft Research Technical Report}, 2025.

\bibitem{Young2025trust}
H.~Young,
``The Trust Algebra: Graded Evidence in Judgment Geometry,''
\emph{Microsoft Research Technical Report}, 2025.

\bibitem{Young2025descent}
H.~Young,
``Descent Obstructions in Program Verification,''
\emph{Microsoft Research Technical Report}, 2025.

\bibitem{Young2025smt}
H.~Young,
``SMT Dispatch: Fragment Routing for Judgment Geometry,''
\emph{Microsoft Research Technical Report}, 2025.

\bibitem{Young2025pcp}
H.~Young,
``Verification Certificates That Ship With Code,''
\emph{Microsoft Research Technical Report}, 2025.

\bibitem{Young2025chains}
H.~Young,
``Certificate Chains for Compositional Verification,''
\emph{Microsoft Research Technical Report}, 2025.

\bibitem{Chen2021Codex}
M.~Chen \emph{et al.},
``Evaluating Large Language Models Trained on Code,''
\emph{arXiv preprint arXiv:2107.03374}, 2021.

\bibitem{Li2023StarCoder}
R.~Li \emph{et al.},
``StarCoder: May the Source Be with You!,''
\emph{arXiv preprint arXiv:2305.06161}, 2023.

\bibitem{Roziere2023CodeLlama}
B.~Rozi\`{e}re \emph{et al.},
``Code Llama: Open Foundation Models for Code,''
\emph{arXiv preprint arXiv:2308.12950}, 2023.

\bibitem{Li2022AlphaCode}
Y.~Li \emph{et al.},
``Competition-Level Code Generation with AlphaCode,''
\emph{Science}, vol.~378, no.~6624, 2022.

\bibitem{Sun2024Clover}
C.~Sun \emph{et al.},
``Clover: Closed-Loop Verifiable Code Generation,''
\emph{arXiv preprint arXiv:2310.17807}, 2024.

\bibitem{First2023Baldur}
E.~First \emph{et al.},
``Baldur: Whole-Proof Generation and Repair with Large Language Models,''
\emph{Proc.\ FSE}, 2023.

\bibitem{Wang2024LEGOProver}
H.~Wang \emph{et al.},
``LEGO-Prover: Neural Theorem Proving with Growing Libraries of Lemmas,''
\emph{Proc.\ ICLR}, 2024.

\bibitem{Necula1997PCC}
G.~C.~Necula,
``Proof-Carrying Code,''
\emph{Proc.\ POPL}, pp.~106--119, 1997.

\bibitem{Ernst2007Daikon}
M.~D.~Ernst \emph{et al.},
``The Daikon System for Dynamic Detection of Likely Invariants,''
\emph{Science of Computer Programming}, vol.~69, no.~1--3, 2007.

\bibitem{Molina2022SpecFuzzer}
F.~Molina \emph{et al.},
``SpecFuzzer: A Tool for Mining Specifications via Fuzzing,''
\emph{Proc.\ ASE}, 2022.

\bibitem{Goguen1992Sheaves}
J.~A.~Goguen,
``Sheaf Semantics for Concurrent Interacting Objects,''
\emph{Mathematical Structures in Computer Science}, vol.~2, no.~2, 1992.

\bibitem{Robinson2014Topological}
M.~Robinson,
\emph{Topological Signal Processing},
Springer, 2014.

\bibitem{Curry2014Sheaves}
J.~Curry,
``Sheaves, Cosheaves and Applications,''
\emph{arXiv preprint arXiv:1303.3255}, 2014.

\bibitem{Hansen2019Toward}
J.~Hansen and R.~Ghrist,
``Toward a Spectral Theory of Cellular Sheaves,''
\emph{Journal of Applied and Computational Topology}, vol.~3, 2019.

\bibitem{Lamsweerde2009Requirements}
A.~van Lamsweerde,
\emph{Requirements Engineering: From System Goals to UML Models to
Software Specifications},
Wiley, 2009.

\bibitem{Gervasi2005Reasoning}
V.~Gervasi and D.~Zowghi,
``Reasoning about Inconsistencies in Natural Language Requirements,''
\emph{ACM TOSEM}, vol.~14, no.~3, 2005.

\bibitem{Chen2021}
M.~Chen \emph{et al.},
``Evaluating Large Language Models Trained on Code,''
\emph{arXiv preprint arXiv:2107.03374}, 2021.

\bibitem{Austin2021MBPP}
J.~Austin \emph{et al.},
``Program Synthesis with Large Language Models,''
\emph{arXiv preprint arXiv:2108.07732}, 2021.

\bibitem{DeMoura2008Z3}
L.~de~Moura and N.~Bj{\o}rner,
``Z3: An Efficient SMT Solver,''
\emph{Proc.\ TACAS}, LNCS vol.~4963, pp.~337--340, 2008.

\end{thebibliography}

\end{document}
